% Standalone section fragment; requires amsmath, amssymb, amsthm, longtable.
\section{The coefficient-support lemma and an explicit recomputation of LMR}
\subsection{S-001. The coefficient-support lemma, including every copy}
Let $n,d\geq1$, and let $E$ be a complex vector space with ordered basis $e_1,\ldots,e_r$.
Use the upper-triangular highest-weight convention, so $E_{i,i+1}$
sends $e_{i+1}$ to $e_i$.  A polynomial in
$\operatorname{Sym}^d(\operatorname{Sym}^n E)$ is written in the algebraically
independent generators $c_\alpha=e_1^{\alpha_1}\cdots e_r^{\alpha_r}$,
$|\alpha|=n$.

\paragraph{Statement.}
If $P$ is any highest-weight vector of weight $\lambda\vdash nd$, every
generator $c_\alpha$ in every monomial occurring with nonzero coefficient
in $P$ satisfies
\[
             \sum_{i=j}^r\alpha_i\leq\lambda_j
                   \qquad(1\leq j\leq r).
\]
For $j=1$ the inequality is automatic whenever the highest-weight space
is nonzero. In particular the statement requested as F-016 holds for
all $j\geq2$, for the complete multiplicity space.

\paragraph{Source correction.}
The attached latest GCT\_KRON paper already contains a numbered finding
F-016 and this support argument: \texttt{GCT\_KRON\_work/GCT\_KRON.tex},
subsection ``Finding F-016: tail-sum support and full determinant
multiplicity'', equivalently
\texttt{agents/formulas/nested\_support.tex}, first proof.
The intake's statement that the numbered full proof is absent is therefore
outdated for the supplied checkpoint. The following is an independent
complete proof of the requested assertion, not a claim that the lemma was
first proved in this continuation.

\begin{proof}
Introduce the matrix-coordinate algebra
$A=\operatorname{Sym}(E\otimes\mathbb C^d)$, with generators $x_{ih}$,
$1\leq i\leq r$, $1\leq h\leq d$. Its subspace of column multidegree
$(n,\ldots,n)$ is canonically $(\operatorname{Sym}^nE)^{\otimes d}$.
Symmetrizing the columns embeds $\operatorname{Sym}^d(\operatorname{Sym}^nE)$
in this subspace, equivariantly for $GL(E)$. The image of a monomial in
the $c_\alpha$ is a nonzero scalar times the sum of its distinct column
assignments. Two distinct coefficient monomials have disjoint sets of
column assignments.

By the Cauchy decomposition
\[
 A_{nd}=\bigoplus_{\mu\vdash nd}S_\mu E\otimes S_\mu\mathbb C^d,
\]
the row highest-weight space of weight $\lambda$ is one copy of the
irreducible column module $S_\lambda\mathbb C^d$. If this module is zero,
there is no assertion to prove. Otherwise $\ell(\lambda)\leq d$.
For $1\leq h\leq\ell(\lambda)$ let $\Delta_h$ be the minor in the
first $h$ rows and first $h$ columns. The nonzero polynomial
\[
                 \prod_{h\geq1}\Delta_h^{\lambda_h-\lambda_{h+1}}
\]
is row highest-weight of weight $\lambda$. Its $GL_d$ column span is the
whole row highest-weight space by irreducibility. Under a column change,
Cauchy--Binet expresses a height-$h$ flag minor as a linear combination
of minors with the same first $h$ rows. Consequently the entire row
highest-weight space is spanned by products of flag minors whose heights
are the column lengths of $\lambda$.

Exactly $\lambda_j$ factors of each spanning product have height at
least $j$. A minor of smaller height has no variable from rows
$j,\ldots,r$. A determinant monomial in any other minor has at most one
variable in a fixed matrix column. Thus every matrix monomial in every
spanning product has degree at most $\lambda_j$ in the entries of rows
$j,\ldots,r$ of each fixed column. Linear combination preserves this
support property. Projecting onto column multidegree $(n,\ldots,n)$ also
preserves it, since that projection only deletes monomials.

Apply this bound to the symmetrization of $P$. If one coefficient
monomial of $P$ contained a forbidden generator, every column assignment
containing that generator would violate the matrix-column bound. Its
coefficient in the symmetrization is nonzero and cannot cancel with a
different coefficient monomial, by the disjoint-support observation.
This contradiction proves the bound for every monomial of $P$.
The argument applies to an arbitrary vector in the full highest-weight
space, so it does not choose or discard multiplicity copies.
\end{proof}

\paragraph{Dictionary and scope.}
This proof is a calculation with exterior powers of a universal map of
locally free sheaves and its torus weight grading. Its representation
ring shadow is the Cauchy identity; that shadow alone does not establish
coefficient support or the rank of a coherent quotient. In the supplied
Grothendieck archive transcription, cote~133, scan~028, the Cauchy tensor
power identity and the question of an actual direct summand are the
relevant exact passage. The affine group-scheme representation/comodule
and diagonalizable-group grading dictionary is SGA~3, Expos\'e~I,
Propositions~4.7.2--4.7.3. The use of irreducibility and symmetrization
above is on the characteristic-zero fiber. No positive-characteristic
semisimplicity, purity, or motivic assertion is inferred.

\subsection{S-002. Six explicit degree-twelve highest-weight functions}
Let $V=\mathbb C^7$ with basis $e_0,\ldots,e_6$, and write a cubic
$F(v)=T(v,v,v)$ with symmetric trilinear form $T$. Put
\[
 a=T_{000},\quad b=T_{001},\quad c=T_{011},\quad d=T_{111},
 \qquad S_i=(T_{ijk})_{0\leq j,k\leq6}\quad(i=0,1),
 \qquad A=(T_{00j})_{0\leq j\leq6}.
\]
Here $d$ in this subsection is a binary cubic coefficient; the outer
polynomial degree is always twelve. Define
\[
 \begin{split}
 H(X,Y)&=\det(XS_0+YS_1)
       =\sum_{k=0}^7\binom7k h_kX^{7-k}Y^k,\\
 B(X,Y)&=\det\begin{pmatrix}XS_0+YS_1&A\\A^t&0\end{pmatrix}
       =\sum_{k=0}^6\binom6k b_kX^{6-k}Y^k.
 \end{split}
\]
The lower-case $b$ without a subscript is $T_{001}$; $b_k$ denotes a
bordered-determinant coefficient. The following six explicit functions
have degree twelve in $T$:
\[
\begin{aligned}
 P_0={}&h_0(2b^5-5ab^3c+3a^2bc^2+a^2b^2d-a^3cd),\\
 P_1={}&h_0(b^5-2ab^3c+a^2bc^2)
        +h_1(-ab^4+2a^2b^2c-a^3c^2),\\
 P_2={}&h_0(2b^5-3ab^3c+a^2b^2d)
        +h_1(-4ab^4+6a^2b^2c-2a^3bd)\\
       &+h_2(2a^2b^3-3a^3bc+a^4d),\\
 P_3={}&h_0(b^5-ab^3c)+h_1(-3ab^4+3a^2b^2c)\\
       &+h_2(3a^2b^3-3a^3bc)+h_3(-a^3b^2+a^4c),\\
 P_4={}&h_0b^5-5h_1ab^4+10h_2a^2b^3-10h_3a^3b^2
        +5h_4a^4b-h_5a^5,\\
 R={}&b_0(5b^3c-3abc^2-ab^2d+a^2cd)-10b_1b^4
       +20b_2ab^3-20b_3a^2b^2+10b_4a^3b-2b_5a^4.
\end{aligned}
\]
They are highest-weight vectors of weight
\[
                         \lambda=(19,7,2,2,2,2,2).
\]
They form a basis of the entire ambient highest-weight space. In
particular $R$ supplies a copy outside the degree-twelve binary
cubic--Hessian-determinant span of $P_0,\ldots,P_4$.

\begin{proof}[Highest-weight calculation]
Let $g$ be upper triangular. Polynomial functions in
$\operatorname{Sym}^{12}(\operatorname{Sym}^3V)$ are evaluated with the
convention $(gP)(F)=P(F\circ g)$. If $g$ fixes $e_0,e_1$ pointwise,
then $S_i(F\circ g)=g^tS_i(F)g$ and $A(F\circ g)=g^tA(F)$.
Both determinant constructions therefore acquire the common factor
$\det(g)^2$. Thus the unipotent radical fixing $e_0,e_1$ fixes all six
functions.

For the remaining simple raising operator $D=E_{0,1}$, its action is
\[
 Da=0,\quad Db=a,\quad Dc=2b,\quad Dd=3c,
 \qquad Dh_k=kh_{k-1},\qquad Db_k=kb_{k-1}.
\]
For the last identity $A$ is fixed because $e_0$ is fixed; the coefficient
normalizations by $\binom7k$ and $\binom6k$ give exactly the displayed
multipliers. Substitution and the product rule give $DP_i=DR=0$.
The standalone checker verifies these six identities by exact sparse
integer arithmetic, with no evaluations or probabilistic step.

Under a diagonal $g=\operatorname{diag}(t_0,\ldots,t_6)$,
$h_k$ has character $\det(g)^2t_0^{7-k}t_1^k$ and $b_k$ has character
$\det(g)^2t_0^{10-k}t_1^k$. Each displayed summand therefore has character
$\det(g)^2t_0^{17}t_1^5$, which is the claimed partition. The $h_k$ have
coefficient degree seven and the $b_k$ degree eight, proving degree twelve.
This establishes the highest-weight assertion directly, without
coefficient-support assumptions.
\end{proof}

\begin{proof}[Ambient dimension and independence]
The ambient multiplicity is independently evaluated over $\mathbb Q$ by
\texttt{check\_lmr\_ambient.py}. Its complete mathematical prescription is
\[
 h_0[h_3]=1,\qquad
 k h_k[h_3]=\sum_{i=1}^k h_{k-i}[h_3]
          \frac{p_i^3+3p_i p_{2i}+2p_{3i}}6,
 \qquad
 a_\lambda=\langle h_{12}[h_3],s_\lambda\rangle.
\]
The code expands the recurrence in power sums using exact rational
arithmetic and computes $\langle p_\rho,s_\lambda\rangle=\chi^\lambda(\rho)$
by Murnaghan--Nakayama beta-number moves. There are 3560 nonzero
power-sum types; the result is exactly six. Independently of the final
value, the cycle-index mass is checked to equal one and the identity
character is checked against the hook-length value
$919850631932232$.

The six displayed functions have a nonzero evaluation determinant
$921522$ modulo $1000003$ on the six generic integral symmetric tensors
specified by the seed rule in S-003 below. Hence they are linearly
independent over $\mathbb Q$. As the ambient multiplicity is six, they
form a basis. This finite-field nonzero minor proves characteristic-zero
independence: every denominator in the definitions is invertible at the
chosen prime, so a zero rational minor would reduce to zero there.
No zero modular minor is used to prove an upper rank bound.
\end{proof}

\subsection{S-003. The exact multiplicities $a=6$, $K_3=5$, $q=6$}
For the partition and degree of S-002, extending the partition by two
zero parts to $GL_9$ (plethysm multiplicities are stable once the ambient
vector-space rank contains the partition), one has
\[
 a_\lambda=6,\qquad
 [S_\lambda\mathbb C^9:I(\operatorname{Det}_3)_{12}]=1,
 \qquad K_3(\lambda)=5,\qquad q_\lambda=6.
\]
An explicit highest-weight equation separating the permanent and
 determinant closures is
\[
                  \Psi=2P_0+21P_1+7P_2+35P_3+7P_4.
\]
This verifies independently the two numerical assertions in
Landsberg--Manivel--Ressayre, \emph{Hypersurfaces with degenerate duals
and the Geometric Complexity Theory Program}, Comment.
Math.
Helv.
88 (2013), \S3.2, p.~476, and recomputes both sides of GCT\_KRON F-025.
The original numbers are published results; no priority claim is made.

\begin{proof}[The equation vanishes on the determinant closure]
Write $f(X,Y)=aX^3+3bX^2Y+3cXY^2+dY^3$.
If $H=fQ$ for a binary quartic
$Q=\sum_{j=0}^4\binom4j q_jX^{4-j}Y^j$, its normalized coefficients obey
\[
 h_k=\binom7k^{-1}\sum_{i+j=k}\binom3i\binom4j f_iq_j,
 \qquad(f_0,f_1,f_2,f_3)=(a,b,c,d).
\]
Substituting these expressions in $\Psi$ gives the zero polynomial,
coefficient by coefficient over $\mathbb Q$. This identity is checked
exactly by \texttt{verify\_lmr\_explicit.py}; the identity also follows by
multiplying out the six displayed formulas.

Now let $F(x)=\det_3(\sum_{i=0}^6x_iM_i)$. At a matrix of rank two,
the Hessian of $\det_3$ has rank six. Indeed left and right invertible
changes reduce it to $\operatorname{diag}(1,1,0)$, where the quadratic
term is
\[
              x_{33}(x_{11}+x_{22})-x_{13}x_{31}-x_{23}x_{32},
\]
a quadratic form of rank six. Here the Hessian of $F$ at $Xe_0+Ye_1$ is $6(XS_0+YS_1)$,
so its determinant differs from $H$ only by the nonzero constant $6^7$.
Consequently the Hessian pulled back to
any seven-dimensional linear parameter space has determinant zero at
such a point. For a dense open set of matrix pencils, the binary cubic
$\det_3(XM_0+YM_1)$ is squarefree and its three roots have matrix rank
two. This open set is nonempty: take $M_0=I$ and $M_1$ diagonal with three
distinct entries. Thus $H$ vanishes at every root of $f$ and $f$ divides
$H$. The preceding polynomial identity gives $\Psi(F)=0$ on that dense
open set, hence on every seven-variable determinant projection by
polynomiality.

Every linear endomorphism substitution of $\det_3$ belongs to its orbit
closure, since $GL_9$ is dense in $\operatorname{End}(\mathbb C^9)$.
The same polynomial evaluated on the first seven variables of any
nine-variable translate therefore vanishes; equivalently $\Psi$ lies
in the full determinant orbit-closure ideal. The construction is
highest-weight for $GL_9$ with two terminal zero weights: the last two
simple raising operators act trivially on functions using only the
first seven coefficient coordinates. This proves the required ideal
inclusion.
\end{proof}

\begin{proof}[Exact evaluation certificates for the entire multiplicity space]
All evaluations below are reconstructed from integers by the checker.
For seed $s$, start $u_0=s$, then
\[
 u_{j+1}=1664525u_j+1013904223\pmod{2^{32}},\qquad
 v_j=\left\lfloor u_{j+1}/2^{16}\right\rfloor\pmod{101}.
\]
For determinant/permanent samples, fill the seven $3\times3$ matrices
$M_0,\ldots,M_6$ in row order from the first 63 values. Use
$F_s(x)=\det_3(\sum x_iM_i)$ or
$F_s(x)=\operatorname{per}_3(\sum x_iM_i)$ as appropriate.
The cubic tensor entries are reconstructed by polarizing the six
permutation terms of the $3\times3$ determinant or permanent and dividing
by six. For generic samples, fill the 84 symmetric tensor entries
$T_{ijk}$, $0\leq i\leq j\leq k\leq6$, in lexicographic order from
the same generator.

The prime is $p=1000003$. The seeds, in order, are
$20261122,\ldots,20261127$. Rows are these samples, and columns are
$P_0,P_1,P_2,P_3,P_4,R$. The complete matrices are in
\texttt{results/lmr\_explicit.json}; the checker rebuilds their entries
from the definitions, not from the stored matrices. It gives:
\[
\begin{array}{c|c|c|r}
 \text{family}&\text{rows}&\text{columns}&\text{minor mod }p\\\hline
 \text{generic cubic}&0,1,2,3,4,5&0,1,2,3,4,5&921522\\
 \det_3&0,1,2,3,4&0,1,2,3,5&170120\\
 \operatorname{per}_3&0,1,2,3,4,5&0,1,2,3,4,5&851602
\end{array}
\]
These nonzero minors imply $K_3(\lambda)\geq5$ and $q_\lambda\geq6$
in characteristic zero. The degree-twelve ambient multiplicity is six
by S-002. On the first permanent sample $\Psi$ has value $715852$
modulo $p$, so it is a nonzero ambient highest-weight vector. Its
identical determinant vanishing therefore gives ideal multiplicity at
least one, and $K_3\leq6-1=5$. Combining the bounds proves
$K_3=5$, ideal multiplicity one, and $q=6$. In particular the entire
six-dimensional isotypic multiplicity space has been compared; the
existence of a single equation is not being mistaken for that comparison.
\end{proof}

\paragraph{Grothendieck dictionary.}
Let $X=\mathbb P(\operatorname{Sym}^3(\mathbb C^9)^*)$ and let
$\mathcal I_D$ be the $GL_9$-equivariant coherent ideal sheaf of the
projective determinant orbit closure. The degree-twelve homogeneous
coordinate sequence, restricted to the $S_\lambda$ multiplicity spaces,
is an exact sequence of dimensions $1\to6\to5$. For the permanent
closure the corresponding kernel has dimension zero and the quotient
has dimension six. The quotient here is the image of ambient homogeneous
sections in the coordinate ring; no projective-normality identification
with all sections of $\mathcal O_D(12)$ is required. These are actual ranks of restriction maps of
coherent sections, not merely relations in $K_0$. The explicit six
functions and their evaluation maps are defined over a localization of
$\mathbb Z$ obtained by inverting six and the binomial denominators.
The nonzero minor modulo $1000003$ proves the asserted generic-fiber
ranks after clearing these denominators. It does not assert flatness or
constant ranks in every characteristic. The exterior determinant
constructions belong to the Schur-operation/Cauchy passage of cote~133,
scan~028, and the equivariant module interpretation is the representation
and comodule dictionary of SGA~3, Expos\'e~I, Proposition~4.7.2. No
$\ell$-adic or motivic statement is needed by, or obtained from, these
calculations.

\subsection{S-004. The independent padding variable and the inherited row bounds}
The current prompt uses $z^{n-m}\operatorname{per}_m$ with $z$ independent
of the $m^2$ matrix variables. Several GCT\_KRON proofs explicitly use
internal padding and the row cap $r\leq m^2$. Under the current convention
the elementary support cap is $r\leq m^2+1$. This change does not
invalidate all the inherited bounds, but it must be checked where used.

\paragraph{Statement.}
In the range $n\geq m^{25}$, the following necessary conditions for an
actual obstruction remain valid for independent padding:
\[
 m\geq190,\qquad r\geq52,\qquad
 \lambda_{2j}\geq28-j\ (1\leq j\leq26),\qquad
 \lambda_3\geq557\ \text{or}\ \lambda_{52}\geq3.
\]
The two-order-ideal construction of M-019, with rectangular root sets,
also proves
\[
              \lambda_4\geq27\quad\text{or}\quad
              \lambda_2\geq238356708.
\]
For $r\leq m^2$ the inherited threshold $239063403$ still applies.
For the extra possible row $r=m^2+1$, the inherited proof does not justify
that stronger numerical threshold. This is a correction of the proof's
scope, not a counterexample to the stronger inequality.

\begin{proof}
F-009 already gives a homogeneous determinantal representation with an
independent auxiliary variable $L$; hence its $m\geq190$ consequence
is unchanged. M-008 represents every $L^{n-m}f$ in a fixed number of
variables and does not use the permanent's row cap. Its $r\geq52$
consequence is likewise unchanged.

For the c=25 staircase in M-014, all cuts other than $k=2$ have $k\geq4$.
Their suffix variable count is $R=r-k+1\leq m^2-2$. The same estimates
in that proof therefore apply with its numerical auxiliary bound $m^2$.
At $k=2$ use M-013 with $\ell=12$, $b=26$ and
$h=\lceil\sqrt{r-1}\rceil\leq m$. Its required size is at most
$2m\binom{m^2+12}{12}$. The ratio of this to $m^{25}$ is
\[
               \frac2{12!}\prod_{i=1}^{12}(1+i/m^2),
\]
which decreases with $m$ and is strictly less than one at $m=190$ by
exact integer comparison. Thus $\lambda_2\leq26$ is excluded and the
same staircase follows.

For A-025 replace the last-group cap $m^2-51$ by $m^2-50$ in its
pruned cover. Explicitly, put
\[
 A(m,v)=(m+1)\binom{49+t}{49}-49\binom{49+t}{50},\qquad t=\min(m,v).
\]
With $u=360,v=195$, the modified count is
\[
 T^{\rm ext}_*(m)=\binom{49+\min(360,m)}{49}
       +A(m,195)+(m^2-50)A(m-1,195).
\]
The standalone checker verifies
$25T^{\rm ext}_*(m)<24m^{25}$ for all 171 integers $190\leq m\leq360$.
For $m\geq360$, write $C=\binom{409}{49}$,
$D=\binom{244}{49}$ and $c=1911/10$. The normalized count is the sum of
\[
 C/m^{25},\qquad D(m+1-c)/m^{25},\qquad
                 D(m^2-50)(m-c)/m^{25}.
\]
Each term is positive and decreasing. For the last term its logarithmic
derivative is at most
$3/m+3/m-25/m<0$; the same inequalities used by A-025 remain valid,
with 50 in place of 51. Thus the checked endpoint establishes the
infinite tail. The same support cover and S-001 then give the same
$557$-or-$3$ necessary condition.

Finally, in M-019 use $m$ inward root trees and $m+1$ outward root trees.
They offer $m(m+1)\geq m^2+1\geq r$ distinct bridges, enough for all
$r-1$ variable labels and the constant label. Every path and cofactor
argument in M-019 is unchanged. With $\ell=12$ each tree has at most
$Q=\binom{m^2+12}{12}$ vertices. If the attached axis ideals both have
$t+1$ members, the required size is at most
\[
                         (2m+1)(t+1)Q.
\]
The function
\[
 \frac{m^{25}}{(2m+1)\binom{m^2+12}{12}}
       =\frac{12!}{(2+1/m)\prod_{i=1}^{12}(1+i/m^2)}
\]
is increasing on $m>0$. At $m=190$ its integer floor is exactly
$J=238356682$, as checked by the two integer inequalities defining a
floor. Set $t=J-1$. If $\lambda_4\leq26$ and
$\lambda_2\leq26+t=238356707$, the support is covered by these two
axis ideals plus the order-26 simplex, exactly as in M-019. The required
matrix size is at most $m^{25}\leq n$, so the whole ambient highest-weight
space survives on the determinant. This forbids an obstruction and
proves the asserted alternative.
\end{proof}

\paragraph{Other consequences of the convention.}
A-035's basic partial-Hessian obstruction exclusion is unaffected:
$m^2+1<n^2$ for every $n>m\geq2$, so the essential dimension remains
proper. Its uniform guarantee that \emph{every} relevant partial-Hessian
binary coefficient algebra is covered by the $r\leq2n$ injection changes
from $m^2\leq2n$ to $m^2+1\leq2n$. The determinant-only support and jet
criteria F-016, M-013, M-015--M-019 are stated in terms of the actual row
count $r$ and retain those exact hypotheses. Internally padded permanent
families used for lower bounds remain available inside the independently
padded closure by specializing the separate padding variable to a matrix
entry, so such lower bounds are not lost. A claim asserting equality or
an upper bound from the smaller essential dimension must be checked
separately, as above.

\subsection{The F-016 dependency inventory}
The complete inventory with exact per-statement proof locators and
subclaim scope is in \texttt{dependencies.tsv}; excluded cross-references
are in \texttt{nondependencies.tsv}. Direct support applications occur
in F-016's own all-copy consequence, F-017, F-018, F-021--F-023,
F-025--F-026, M-011, M-013, M-015--M-019, A-025, A-031,
O-015.b, O-018.a, O-020, O-021 and O-022. The following additional
statements inherit those applications: M-014, O-015.c, O-018.b,
O-023--O-025 and O-027; O-027 also explicitly repeats the support step.
For O-024 the support interpretation is separate from its combinatorial
cover identity. O-026's multiplicity interpretation depends on the
support lemma, whereas its exact combinatorial minima do not.
A-035 uses it only in the optional O-027 full-ambient-coverage corollary,
not in the basic bracket nonvanishing or binary-algebra injection.
O-031, QA033 and QA033-CB use it only for their conditional eventual
multiplicity consequence; their unresolved Jacobian gate is not supplied
by the lemma. M-020 originally inherits it for the illustrative F-025
example, not for its general saturation statement. The specialized
F-025 example now has the independent proof S-003, so that numerical
comparison no longer needs F-016.

Each inventory check confirms the particular inference from coefficient
support to an all-copy ideal or quotient statement. It does not constitute
an independent proof of every other geometric or numerical input to the
referenced finding. In every direct occurrence the same valid step is
used: a highest-weight polynomial belongs to the polynomial algebra on
the indicated suffix-bounded coefficient set. Dominant projection to
that set makes its pullback injective; alternatively an exact determinant
representation of the retained padded forms makes every determinant-ideal
highest-weight vector vanish on the permanent family. Both arguments
apply to all vectors of the multiplicity space. Numerical substitutions
using the row cap must use the corrected convention in S-004.

\subsection{S-005. A finite 804-dimensional gate for the $(8,9)$ question}
This is a proved reduction of one concrete four-row family, not an
explicit obstruction. Let $F$ be a degree-$n$ form in four variables and
write $F(v)=T(v,\ldots,v)$ with symmetric $n$-linear tensor $T$.
Distinguish $e_0$ and put $U=\langle e_1,e_2,e_3\rangle$. For $v\in U$
write
\[
 F_j(v)=T(e_0^{n-j},v^j),\qquad 0\leq j\leq n.
\]
Thus $F_0$ is a scalar. Let $\mathcal H_d(n,c)$ be the ambient
highest-weight space of weight $(nd-3c,c,c,c)$ in outer degree $d$.
Assume $1\leq c\leq n$. Define
\[
 \mathcal R_c=
 \left[\operatorname{Sym}\left(\bigoplus_{j=2}^c
                        \operatorname{Sym}^j\mathbb C^3\right)\right]_{\det^c},
\]
where the bracket denotes the $GL_3$ determinant-character space, so
all its elements have total ordinary $GL_3$ weight $3c$.

\paragraph{Uniform stabilization.}
For every $d\geq3c$, multiplication by $F_0^{d-3c}$ is an isomorphism
\[
             \mathcal H_{3c}(n,c)\longrightarrow\mathcal H_d(n,c),
             \qquad\mathcal H_{3c}(n,c)\simeq\mathcal R_c.
\]
The assertion concerns the actual polynomial highest-weight spaces,
not only equality of their dimensions.

\begin{proof}
Localize at $F_0\ne0$ and define the centered coefficients
\[
 C_j(v)=\sum_{i=0}^j(-1)^{j-i}\binom ji
               F_i(v)\frac{F_1(v)^{j-i}}{F_0^{j-i}}
       =T\left(e_0^{n-j},
          \left(v-\frac{F_1(v)}{F_0}e_0\right)^j\right).
\]
Then $C_1=0$, and $F_0^{j-1}C_j$ is polynomial of coefficient degree $j$.
For the terms $i=0,1$, the apparent denominator $F_0^j$ is reduced by
the factor $F_0$ in $F_i$ or by collecting the two terms; for $i\geq1$
the displayed denominator has degree at most $j-1$ directly.
More simply, the $i=0$ term is $(-1)^jF_1^j/F_0^{j-1}$, and all other
terms have denominator at most $F_0^{j-1}$.

The three-dimensional unipotent group $g(e_i)=e_i+a_ie_0$, $g(e_0)=e_0$,
translates $F_1/F_0$ freely and fixes every $C_j$. The change between
$(F_0,F_1,F_2,\ldots,F_c)$ and
$(F_0,F_1,C_2,\ldots,C_c)$ is triangular and invertible after localizing
at $F_0$. Consequently its invariant ring is exactly
\[
             \mathbb C[F_0,F_0^{-1},C_2,\ldots,C_c].
\]
By S-001, every vector of $\mathcal H_d(n,c)$ uses only the coefficients
$F_0,\ldots,F_c$, since every coefficient generator has tail degree at
most $\lambda_2=c$. The remaining highest-weight condition with tail
$(c,c,c)$ says precisely that its centered expression is $SL_3$ invariant
with $GL_3$ character $\det^c$.

Decompose such an invariant $I$ into pieces $I_k$ with $k$ centered-jet
factors. The weighted degree condition is $\sum j=3c$ in each monomial,
and $j\geq2$, so $k\leq3c/2$. Its denominator is at worst
$F_0^{3c-k}$. Hence
\[
                  L(I)=\sum_kF_0^{3c-k}I_k(C_2,\ldots,C_c)
\]
is a polynomial of coefficient degree $3c$. Under
$\operatorname{diag}(t_0,g)$, $C_j$ transforms as
$t_0^{n-j}\operatorname{Sym}^j g$, so $L(I)$ has character
$t_0^{3cn-3c}\det(g)^c$ and is highest-weight. At $F_0=1,F_1=0$ it
recovers $I$ on freely assignable $C_2,\ldots,C_c$, proving injectivity.

Conversely the localized expression of a degree-$d$ highest-weight
polynomial is $\sum_kF_0^{d-k}I_k(C)$. This follows from coefficient
homogeneity and the tail character, and it has no additional dependence
on the freely translated $F_1/F_0$. It equals $F_0^{d-3c}L(I)$.
For $d\geq3c$ this is a polynomial identity; no division by $F_0$ is
being asserted in the unlocalized coordinate ring. The same argument
shows that multiplication by $F_0^{3c-d}$ embeds every earlier
$\mathcal H_d(n,c)$ into $\mathcal H_{3c}(n,c)$. This proves all claims.
\end{proof}

\paragraph{The concrete dimension.}
At $n=9,c=8$ one obtains
\[
          \dim\mathcal H_{24}(9,8)
                  =a_{(192,8,8,8)}=804.
\]
The standard-library checker \texttt{check\_centered\_jet\_gate.py}
performs two independent character computations. The centered-jet
computation takes the coefficient of the determinant character using
six $GL_3$ Weyl terms in
\[
          \prod_{\substack{\alpha\in\mathbb Z_{\geq0}^3\\
                           2\leq|\alpha|\leq8}}(1-x^\alpha)^{-1}.
\]
Its signed dimensions are
\[
       2604128-2468696-2468696+2199430+2237895-2103257=804.
\]
The second computation applies all 24 $GL_4$ Weyl terms directly to
$\operatorname{Sym}^{24}(\operatorname{Sym}^9\mathbb C^4)$ at
$(192,8,8,8)$ and gives the same value. In these terms the total tail
weight is at most 24; each monomial in nonconstant tail generators
therefore uses at most 24 factors and has a unique completion to outer
degree 24 by powers of the pure first-coordinate generator. This
justifies the checker's three-dimensional integer knapsack for the
direct plethysm calculation.

\paragraph{The padded side is completely known in this family.}
For $(m,n)=(8,9)$, CUE\_GCT R-13/R-14 identify the four-variable padded
image with all forms $\ell q$, where $q$ is an arbitrary octic.
It contains the complete linear subspace $e_0^*\operatorname{Sym}^8V^*$.
Its coefficients $F_0,\ldots,F_8$ vary freely, and S-001 puts every
vector of $\mathcal H_d(9,8)$ in their polynomial algebra. Restriction
to this subspace is therefore injective on the entire highest-weight
space, proving
\[
                  q_{(9d-24,8,8,8)}=a_{(9d-24,8,8,8)}
\]
whenever the displayed tuple is a partition. In particular $q=804$
at $d=24$.

\paragraph{One finite rank decides all degrees of this family.}
The determinant coordinate ring is a domain and $F_0$ is nonzero on it.
Thus an ideal vector in degree $d>24$ factors as
$F_0^{d-24}$ times one in degree 24 by the proved polynomial identity,
and that factor is again in the determinant ideal. A nonzero ideal
vector in an earlier degree can be multiplied by $F_0^{24-d}$.
Therefore a gap exists in this family in \emph{some} degree if and only
if the restriction map from the explicit 804-dimensional space
$\mathcal H_{24}(9,8)$ to the determinant coordinate ring has rank below
804. This reduction resolves the all-degree issue without asserting any
uncomputed determinant-side rank.

\paragraph{A rank-803-or-804 dichotomy, using this run's jet certificate.}
Let $Y$ be the closure of the determinant image on the 165 coefficients
$F_0,\ldots,F_8$. The root certificate
\texttt{root/results/det9\_nested\_profiles.json}, profile $(8,8,7)$,
has rank 164 and minor 76 modulo 101. Combined with the determinant
dimension upper bound 164, it makes $Y$ an irreducible hypersurface.
The same full-rank subprojection contains every coefficient through
order seven, proving dominance of that smaller jet.

Since the line-stabilizer parabolic preserves the 165-coordinate
algebra, it preserves the principal prime ideal of $Y$. Its irreducible
homogeneous generator $G$ is therefore a parabolic semi-invariant of
weight $(9D-3c_G,c_G,c_G,c_G)$ for some nonnegative integer $c_G$.
If $c_G\leq7$, S-001 places $G$ in the dominant order-seven jet algebra,
a contradiction. Hence $c_G\geq8$.

If a nonzero vector of $\mathcal H_{24}(9,8)$ is divisible by $G$, its
polynomial quotient is a highest-weight semi-invariant of tail weight
$(8-c_G,8-c_G,8-c_G)$. Negative tail weights are impossible in the
polynomial algebra, so $c_G>8$ gives no kernel. If $c_G=8$, the quotient
has tail weight zero and is a power of $F_0$ times a scalar. Moreover
$G$ has coefficient degree $D\leq24$: otherwise the stabilization
identity would make $G$ divisible by $F_0$, contradicting irreducibility
and the fact that $Y$ is not the hyperplane $F_0=0$. Thus the kernel
has dimension exactly one when $c_G=8$, and zero when $c_G>8$.
Consequently the two possible outcomes of the finite calculation are
\[
  q=804>K_9=803\quad(c_G=8),\qquad
  q=K_9=804\quad(c_G>8).
\]
Finding S-006 below computes the finite rank as 804 and therefore
decides the second alternative.  In particular $c_G\geq9$.  The full
value of $c_G$ and the defining polynomial $G$ remain unknown.

\paragraph{Geometric dictionary.}
Centering is an explicit coordinate calculation on the principal affine
open $D(F_0)$ and a quotient by a split unipotent group. Clearing powers
of $F_0$ constructs actual homogeneous sections. The final rank test is
the $S_{(192,8,8,8)}$ part of the equivariant coherent ideal of the
quaternary determinant image; the principal ideal $(G)$ describes its
projection to the octic-polar coefficient scheme. The all-degree
reduction uses localization and the fact that multiplication by a
nonzero section on an integral scheme is injective. It does not replace
the coordinate ring by its normalization or infer ranks from its
$K_0$ class. The representation/comodule interpretation is again SGA~3,
Expos\'e~I, Proposition~4.7.2, with Schur operations as in the supplied
cote~133, scan~028. No assertion about weights of $\ell$-adic cohomology
or motives is needed by this affine calculation.

\subsection*{S-006. An explicit 804-function certificate excludes the entire equal-tail-eight family}

\paragraph{Statement.}
For every integer $d\geq4$, put
\[
                 \lambda_d=(9d-24,8,8,8).
\]
For the independent padding convention of the present prompt,
\[
                 K_9(\lambda_d)=q_{\lambda_d}=a_{\lambda_d}.
\]
These common multiplicities are 804 for $d\geq24$.
Thus this entire all-degree family produces no multiplicity obstruction
for $(m,n)=(8,9)$.  The irreducible octic-polar determinant hypersurface
of S-005 has tail character $c_G\geq9$.

\paragraph{Explicit invariant polynomials.}
For each profile $\boldsymbol{k}=(k_2,\ldots,k_8)$ satisfying
$\sum_{j=2}^8jk_j=24$, introduce $k_j$ vertices of degree $j$.
An eight-bracket graph is a list of eight ordered triples of distinct
vertices, with each vertex occurring as often as its degree.
If the triples are $e=1,\ldots,8$, write $i_{e,s}\in\{0,1,2\}$ for
the three colors on triple $e$, and let $\epsilon_{i_{e,1}i_{e,2}i_{e,3}}$
be the alternating symbol.  The graph defines the completely explicit
integer polynomial
\[
 I_\Gamma(C_2,\ldots,C_8)=
 \sum_{(i_{e,s})}
 \left(\prod_{e=1}^8\epsilon_{i_{e,1}i_{e,2}i_{e,3}}\right)
 \prod_v C_{j_v}\bigl(i_{e,s}:(e,s)\text{ is incident to }v\bigr).
\]
The symmetric tensor entries are the normalized polarization entries,
so no multinomial coefficient is inserted into this contraction.
Each alternating triple transforms by the determinant character.
Therefore $I_\Gamma$ transforms under $\mathrm{GL}_3$ by $\det^8$,
has torus weight $(8,8,8)$, and is killed by both simple raising
operators.  This also follows by expanding the alternating-symbol
identity, without invoking a spanning theorem for tensor invariants.

The 804 graphs and all their expanded integral coefficients are supplied
in \texttt{support\_lmr/results/bracket\_basis/}.  The profile files
record their graph lists, coefficient tables, and pivot monomials.
The independent multigraded character computation
\texttt{invariant\_profiles.py} gives 154 nonzero profiles and total
dimension 804; the largest profile has dimension 29.  In every profile,
the selected coefficient minor is nonzero modulo 101.  Thus the stored
804 integer polynomials are linearly independent over $\mathbb Q$.
S-005's independent total character calculation bounds the whole
invariant space by 804, so they form a complete basis.

\paragraph{Compression and an exact coefficient check.}
The stored coefficient tables use a lossless sixfold color-symmetry
compression.  Fix the first bracket to the color order $(0,1,2)$ and
expand the remaining seven brackets.  Let $S_O$ be the sum of its
coefficients over a global color-permutation orbit $O$ of monomials.
The full coefficient of each monomial of $O$ is
\[
                             \frac6{|O|}S_O.
\]
Indeed, restoring the first bracket sums over six global color
permutations; each contributes sign $(\operatorname{sgn}\sigma)^8=1$,
and the stabilizer of an orbit element has order $6/|O|$.
In particular the compressed coefficient rank differs from the full
coefficient rank only by nonzero column scales modulo 101.

The standalone standard-library checker
\texttt{verify\_bracket\_basis.py} reconstructs these integer
polynomials and checks their raising identity exactly, rather than
trusting the graph-expansion program.  In coordinates
$T_{a,b,c}$, the first raising operator is
\[
                  E_{01}T_{a,b,c}=bT_{a+1,b-1,c}.
\]
All collected coefficients of $E_{01}I_\Gamma$ are zero as integers.
Global color symmetry conjugates this identity to that for $E_{12}$.
The checker also verifies weight $(8,8,8)$ and every profile minor.
It reads the saved NumPy array format using only
\texttt{zipfile}, \texttt{struct}, and \texttt{ast}; NumPy is not
required by this verifier.

\paragraph{Lifting to the required highest-weight space.}
For a graph with $k=\sum_jk_j$ vertices, define
\[
                    L(I_\Gamma)=F_0^{24-k}I_\Gamma(C_2,\ldots,C_8),
\]
where the centered tensors are exactly those of S-005.
Each $C_j$ has denominator dividing $F_0^{j-1}$, so the displayed
expression is a polynomial.  It has coefficient degree 24 and
highest weight $(192,8,8,8)$.  S-005 proves that these 804 lifts are a
basis of the complete highest-weight multiplicity space, not merely a
collection of copies with the desired weight.

\paragraph{Genuine determinant evaluations.}
Take
\[
     F=x_0^9\det\left(I_9+
              (x_1/x_0)A_1+(x_2/x_0)A_2+(x_3/x_0)A_3\right),
\]
where the three integer matrices have exactly zero traces.  Then
$F_0=1$, $F_1=0$, and $C_j=F_j$.  If $f_{a,b,c}$ is the ordinary
coefficient of $x_0^{9-j}x_1^ax_2^bx_3^c$, where $a+b+c=j$, its
normalized tensor entry is
\[
             T_{a,b,c}=\frac{(9-j)!a!b!c!}{9!}f_{a,b,c}.
\]
All denominators are invertible modulo 101.  The deterministic sampler
\texttt{quaternary/scripts/centered\_det9\_sampler.py} generates 830
such pencils from the specified seed 20261122, reconstructs every
coefficient by exact total-degree Newton interpolation, and checks
direct determinant values.  Its proof and independent reconstruction
checks are recorded in QC-007.  The final negative diagonal entry
makes each source trace zero over $\mathbb Z$, not only in the residue
field, so every sample lifts to a centered rational determinant pencil.

Evaluation of the complete 804-function basis gives an
$804\times830$ matrix over $\mathbb F_{101}$ with rank 804.
The selected maximal minor is 70 modulo 101.  Its serialized row-major
byte matrix has SHA-256
\begin{center}\small\ttfamily
71b43abba722f968438ea9a4c77ad7f899f5a1dd6b8666962bdceabc8fd50d2d.
\end{center}
The matrix, profile ordering, selected columns, and minor are supplied
in \texttt{support\_lmr/results/bracket\_evaluation/} and
\texttt{bracket\_rank\_packed.json}.  The complete standard-library
entry verifier and the independent scalar rank checker are supplied
by QC-007.  All arithmetic bounds for their packed operations are
explicitly checked; the scalar Gaussian calculation is independent
of packed elimination.

A nonzero minor modulo 101 proves a nonzero rational minor at the
lifted determinant pencils.  Hence all 804 highest-weight functions
survive independently on the determinant orbit closure.  The ambient
upper bound and padded equality of S-005 give
$K_9=q=a=804$ at degree 24.

\paragraph{All-degree conclusion.}
S-005 proves that multiplication by $F_0^{d-24}$ identifies the entire
highest-weight spaces for $d\geq24$, and that a nonzero ideal vector
in an earlier admissible degree would remain nonzero after multiplication
to degree 24.  The determinant coordinate ring is integral and $F_0$
is nonzero there.  Thus the degree-24 zero kernel excludes ideal vectors
in every $\lambda_d$ of the family.  The padded restriction is injective
on the same entire spaces.  This proves the statement for every
$d\geq4$, and the hypersurface conclusion follows from S-005's
$c_G=8$ versus $c_G>8$ dichotomy.

\paragraph{The first stable outer degree is sixteen.}
The additional standalone checker
\texttt{check\_equal\_tail\_stabilization.py} counts monomials by
their three tail weights and number of non-pure-first-coordinate
factors, then applies the full 24-term Weyl denominator.  For
$\lambda_d=(9d-24,8,8,8)$, the multiplicities in degrees $4$ through
$16$ are
\[
       1,5,26,89,214,377,543,670,750,787,800,803,804.
\]
Multiplication by $F_0$ is injective, while the complete centered
space has dimension 804.  Thus the common multiplicity is already
804 for every $d\geq16$.  This strengthens the sufficient stable
degree 24 used to write the explicit polynomial lifts; it does not
change their validity.
For the next equal-tail-nine family the same exact computation gives
\[
       0,0,13,92,324,736,1248,1717,2066,2269,2358,2383,2386
\]
in degrees $4$ through $16$, and the complete centered upper bound
is 2386.  Its ambient multiplicity likewise stabilizes first at
degree 16.  S-007 proves its padded multiplicity is ambient, but
its determinant restriction still requires a separate rank calculation.

\paragraph{A failed smaller evaluation family, exactly quantified.}
Diagonal pencils produce products of nine linear forms.  Their
degree-24 coordinate-ring multiplicity is bounded by the multiplicity
of $S_{(192,8,8,8)}$ in
$\operatorname{Sym}^9(\operatorname{Sym}^{24}\mathbb C^4)$, which is
377.  Thus they could never certify rank 804.  In contrast, the
normalization target for products of three arbitrary cubics is
$\operatorname{Sym}^3(\operatorname{Sym}^{24}
(\operatorname{Sym}^3\mathbb C^4))$, with multiplicity 699257 in this
type, so this dimension cap gives no obstruction to using that family.
Both exact 24-term Weyl calculations are recomputed by the standalone
\texttt{check\_small\_family\_caps.py}.  These are upper bounds from
parameter-space pullback, not assertions that the coordinate rings
equal their normalization targets.

\paragraph{Geometric dictionary and scope.}
The certificate computes the $S_{(192,8,8,8)}$ isotypic part of the
equivariant homogeneous restriction morphism to the integral determinant
scheme.  Its kernel is zero; localization at the nonzero section $F_0$
and the explicit denominator clearing propagate this fact through the
whole equal-tail-eight family.  The Schur operation lives in the
representation category and its $K_0$, while the conclusion concerns
the actual coherent ideal and its graded sections.  The dictionary is
SGA~3, Expos\'e~I, Proposition~4.7.2, together with the supplied
cote~133, scan~028 Cauchy/Schur-functor discussion.  The modular minor
proves the characteristic-zero statement only; neither an
all-characteristic assertion nor an $\ell$-adic or motivic conclusion
is inferred.

\subsection{S-007. An extra column of full padded multiplicity}
\paragraph{Statement and conventions.}
Work over $\mathbb C$, let $m\geq2$ be even, and let $n>m$.
Let $V_{m,n}$ be the closure of the quaternary forms
$\ell^{n-m}q_m$, where both the linear form $\ell$ and the degree-$m$
form $q_m$ vary freely.  For every four-row partition
$\lambda\vdash nd$ with $\lambda_2\leq m+1$, restriction to this
variety is injective on the complete highest-weight space:
\[
 \operatorname{mult}_{S_\lambda}\mathbb C[V_{m,n}]_d
       =a_\lambda(d[n]).                                      \tag{S-007.1}
\]
There is no lower bound on the outer degree beyond the requirement
that $\lambda$ be a partition.  In particular, for the actual padded
permanent closures in the question,
\[
 \begin{array}{ll}
 (m,n)=(8,9):&q_\lambda=a_\lambda\quad(\lambda_2\leq9),\\
 (m,n)=(17,23):&q_\lambda=a_\lambda\quad(\lambda_2\leq17).
 \end{array}                                                  \tag{S-007.2}
\]
The second assertion uses the checked dominance of the quaternary
degree-16 permanent image and the block specialization
$\operatorname{per}_{17}\rightsquigarrow z\operatorname{per}_{16}$.
The first uses the checked degree-eight dominance.  These are the
permanent-image inputs supplied in CUE\_GCT and retained by OBS-004;
the proof below concerns the padding step and every multiplicity copy.

\paragraph{Centering preserves the whole space in every degree.}
Write the normalized coefficients as
\[
 F(x_0,v)=\sum_{j=0}^n\binom nj x_0^{n-j}F_j(v),\qquad
 C_j=\sum_{i=0}^j(-1)^{j-i}\binom ji
                  F_i(F_1/F_0)^{j-i}.
\]
Thus $C_0=F_0$, $C_1=0$, and $C_j$ is homogeneous of coefficient
degree one as an element of the localization at $F_0$.
As proved in S-005, this triangular change of variables identifies
the invariants of the three translations with
$\mathbb C[F_0^{\pm1},C_2,\ldots,C_n]$.
For a highest-weight polynomial of weight
$\lambda=(nd-w,\beta)$, where $\beta=(b,c,e)$ and $w=b+c+e$,
S-001 allows only $F_0,\ldots,F_b$.  Its centered expression therefore
uses only $C_2,\ldots,C_b$.
Specializing $F_0=1$ gives a residual $\mathrm{GL}_3$ highest-weight
polynomial $I$ of weight $\beta$, with weighted degree $w$ when
$C_j$ has weight $j$.
This specialization is injective on the fixed outer-degree space:
if a centered monomial has $k$ factors, its accompanying power of
$F_0$ is uniquely $d-k$.  Distinct such monomials remain distinct
after $F_0=1$.  Negative powers in the intermediate localization do
not affect this argument.  In particular no stabilization assumption
or divisibility assertion about a power of $F_0$ is being suppressed.

\paragraph{The last allowed jet occurs at most linearly.}
Set $s=m+1$, which is odd.  The decomposition
\[
 \operatorname{Sym}^2(\operatorname{Sym}^s\mathbb C^3)
   =\bigoplus_{i=0}^{(s-1)/2}S_{(2s-2i,2i)}\mathbb C^3          \tag{S-007.3}
\]
has first row at least $s+1$ in every summand.  One way to see the
parity in this formula is to decompose the tensor square by Pieri:
the highest vector with $k$ two-letter brackets changes by $(-1)^k$
when the two factors are interchanged, so precisely even $k$ occur
in the symmetric square.  In a tensor product with a polynomial
representation, Littlewood--Richardson containment cannot decrease
the first row.  Higher powers of $C_s$ are quotients of this symmetric
square tensored with further polynomial representations.  Therefore
no sector of $C_s$-degree at least two can contain a type with
$b\leq s$.  Consequently every $I$ under consideration has the form
\[
             I=I_0(C_2,\ldots,C_m)+A(C_2,\ldots,C_m)[C_s],     \tag{S-007.4}
\]
where the bracket denotes linear evaluation in the final tensor.

\paragraph{A centered slice of every padding order.}
Choose $C_2,\ldots,C_m$ and a linear form $l(v)$ arbitrarily, and set
$B_0=1$, $B_1=0$, $B_j=\binom nj C_j$ for $2\leq j\leq m$.
Put $r=n-m$.  In the formal variable $t$, let
\[
 Q_m(t)=\left[(1+lt)^{-r}\sum_{j=0}^m B_jt^j\right]_{\leq m}.
\]
Homogenizing $(1+lt)^rQ_m(t)$ gives a form in $V_{m,n}$ whose
first $m$ jets are the chosen ones and whose first derivative is
zero.  The next ordinary coefficient is
\[
 B_{m+1}=\sum_{j=0}^m(-1)^{m-j}
              \binom{n-j}{m+1-j}l^{m+1-j}B_j.
\]
Indeed it is the negative of the omitted coefficient of degree
$m+1$ in the inverse series.  Using
\[
 \frac{\binom{n-j}{m+1-j}\binom nj}{\binom n{m+1}}
                  =\binom{m+1}j
\]
gives the normalized identity, independent of the padding order,
\[
 C_{m+1}=(-1)^m l^{m+1}
       +\sum_{j=2}^m(-1)^{m-j}\binom{m+1}j l^{m+1-j}C_j.       \tag{S-007.5}
\]
If (S-007.4) vanished on all these forms, its coefficient of highest
degree $m+1$ in $l$ would be $A[l^{m+1}]$.  Pure powers span
$\operatorname{Sym}^{m+1}\mathbb C^3$ in characteristic zero.
Thus $A=0$ for every value of the arbitrary lower jets, and hence
identically; then $I_0=0$.  The injectivity of the centered
specialization proves (S-007.1).  Exactness of highest-weight
decomposition in characteristic zero turns this zero-kernel assertion
into equality of the complete multiplicities.

For $(8,9)$ the full padded variety is available.  For $(17,23)$,
the permanent block specialization gives
$z^6\operatorname{per}_{17}\rightsquigarrow z^7\operatorname{per}_{16}$,
and the checked dense quaternary image of $\operatorname{per}_{16}$
puts $V_{16,23}$ in the actual padded-permanent closure.  Apply
(S-007.1) with even $m=16$.  This proves (S-007.2).
It upgrades the determinant exclusions for tails $(9,6,6)$ and
$(9,7,4)$ to $K=q=a$ in every degree.  At $n=23,d=6$, the newly
covered labels $(87,17,17,17)$, $(88,17,17,16)$,
$(89,17,17,15)$, $(90,17,17,14)$ have ambient multiplicities
$450,1332,2392,3577$, respectively, by OBS-004's extended exact
character calculation.  The two previously checked minimal
first-tail-17 profiles were already covered by OBS-004's projection
certificate; they are not new numerical consequences here.

\paragraph{Scope and geometric dictionary.}
The argument establishes injectivity of the actual graded restriction
map, using a polynomial family in its target scheme and localization
on an invariant affine chart.  It makes no equality assumption about
global sections on a normalization.  The representation/comodule
dictionary is SGA~3, Expos\'e~I, Propositions~4.7.2--4.7.3; the
polynomial Schur-functor and direct-summand interpretation is the
supplied cote~133, scan~028.  The standalone checker
\texttt{check\_even\_padding.py} verifies the inverse-series identity
with exact rational arithmetic for several padding orders; the proof
above establishes it for all $n>m$.

\subsection{S-008. The weight-27 and weight-28 bands with first part at most eleven}
For $(m,n)=(8,9)$, let $\lambda=(9d-27,\beta)$ be a partition,
where $\beta\vdash27$ has at most three parts and $\beta_1\leq11$.
Then
\[
                              q_\lambda=a_\lambda.            \tag{S-008.1}
\]
Together with QC-010, this gives full padded multiplicity for every
tail with $\beta_1\leq11$ and $|\beta|\leq27$.
The proof is an assertion about the full highest-weight space in each
admissible outer degree, using S-007's injective centering map.

\paragraph{The possible cubic coefficients.}
An invariant of weighted degree 27 is at most cubic in $C_9$.
Its coefficients at $C_9$-degrees $0,1,2,3$ have weighted degrees
$27,18,9,0$, respectively, in $C_2,\ldots,C_8$.
The exact six-term Weyl character calculation for
$\operatorname{Sym}^3(\operatorname{Sym}^9\mathbb C^3)$ is
\[
\begin{array}{c|rrrrrrr}
\beta&(9,9,9)&(10,9,8)&(10,10,7)&(11,8,8)&(11,9,7)&(11,10,6)&(11,11,5)\\
\hline
\text{multiplicity}&0&0&1&1&0&0&0.
\end{array}                                                   \tag{S-008.2}
\]
The checker \texttt{check\_cubic\_padding\_boundary.py} enumerates
all unordered triples of the 55 nonic monomials and applies the
six-term Weyl denominator.  The root separately recomputed every
entry with an exact BigInt character dynamic program.

For the five shapes without a cubic component, the quadratic
coefficient has weight nine.  It cannot contain $C_8$, since there
is no generator of weight one.  QC-010's injective map
\[
 \sigma:\operatorname{Sym}^2(\operatorname{Sym}^9E)
       \longrightarrow\operatorname{Sym}^{10}E\otimes\operatorname{Sym}^8E,
 \qquad A\longmapsto A(l^9,lC_8),                              \tag{S-008.3}
\]
then kills that coefficient if the padded restriction is zero:
it is isolated by the terms linear in $C_8$ and of degree ten in
$l$ in (S-007.5).  Each even-bracket summand
$S_{(18-k,k)}$, $k=0,2,4,6,8$, occurs once in the target and maps
nontrivially, which proves the asserted injectivity.  A linear
$C_9$ term has $l$-degree at most nine and cannot interfere.
Once the quadratic part is zero, the affine-linear argument of
S-007 finishes the proof.

\paragraph{Eliminating the two one-dimensional cubic components.}
Use three symbolic nonic letters $a,b,c$, put
$\Delta=[abc]$, and denote the two-coordinate minors by
$E=[ab]$, $F=[ac]$, $G=[bc]$.  Nonzero highest symbols for the
two exceptional components are
\[
 G_1=\Delta^7EFG\quad\bigl(\beta=(10,10,7)\bigr),\qquad
 G_2=\Delta^8a_0b_0c_0\quad\bigl(\beta=(11,8,8)\bigr).       \tag{S-008.4}
\]
Each has degree nine in each letter and is symmetric under the
permutations of $a,b,c$.  For $G_1$, both $\Delta^7$ and $EFG$
change sign under a transposition; for $G_2$ the symmetry is immediate.
The displayed brackets also verify the weights and highest-vector
identities directly.

Specialize all lower jets except $C_7$ to zero.  The coefficient
weights $9,18,27$ of the lower $C_9$-degrees are not multiples of
seven, so all those coefficients vanish.  The padded relation becomes
\[
                              C_9=l^9-36l^2C_7.
\]
The term cubic in $C_7$ is therefore a nonzero scalar times
$G_i(l^2C_7)$.  It remains to check that these restrictions are not
identically zero.

Take $l=e_0$ and in the polarized symbolic polynomials set
$a_2=c_1=0$.  Then
\[
       \Delta=c_2E+a_1b_2c_0,\qquad EFG=a_1b_1c_0^2E.
\]
The coefficient in $G_1$ of
\[
       a_0^4a_1^5\,b_0^2b_1^5b_2^2\,c_0^4c_2^5
\]
is $\binom72\binom62=315$.  The coefficient in $G_2$ of
\[
       a_0^3a_1^6\,b_0^4b_1^2b_2^3\,c_0^4c_2^5
\]
is $-\binom83\binom53=-560$.  In both cases every letter has
its zeroth exponent at least two.  These coefficients therefore
belong to the restriction of the polarized form to
$(e_0^2\operatorname{Sym}^7E)^3$.  Since the polarized form is
symmetric, characteristic-zero polarization implies that its
diagonal cubic restriction to $e_0^2\operatorname{Sym}^7E$ is
nonzero.  The exact integral coefficient calculations are included
in the same standalone checker.

Thus vanishing on the padded slice forces the scalar cubic
coefficient to be zero in both exceptional types.  The remaining
quadratic coefficient again has weight nine and is independent of
$C_8$; (S-008.3) kills it, and the affine-linear argument kills the
rest.  This proves (S-008.1) for all seven shapes, with no assertion
about the determinant-side kernel.

\paragraph{Refinement: the whole weight-28 band also survives.}
For every tail $\beta$ with $|\beta|=28$ and $\beta_1\leq11$,
and every admissible outer degree, one likewise has $q=a$.
The coefficient of $C_9^3$ would have weight one and is therefore
absent.  The quadratic coefficient has weight ten, so its dependence
on $C_8$ is exactly $L(C_2,C_8)$, bilinear in these two tensors.
Since the first row is at most eleven, the quadratic nonic type is
only $S_{(10,8)}$.  Before tensoring with the single quadratic jet
$C_2\in\operatorname{Sym}^2E$, Pieri gives exactly the six types
\[
 (10,8,8),\ (10,9,7),\ (10,10,6),\
 (11,8,7),\ (11,9,6),\ (11,10,5).
\]
The leading map from QC-010, obtained by taking $C_8$-degree two
and $l$-degree ten, is injective on five of these summands and
has kernel exactly the single $S_{(10,9,7)}$.  Tensoring with
$\operatorname{Sym}^2E$ preserves that description of the kernel.

It remains to kill the residual $S_{(10,9,7)}\otimes
\operatorname{Sym}^2E$ coefficient.  Set all lower jets except
$C_2,C_6,C_8$ to zero.  In the padded substitution isolate the
term with respective degrees
\[
                   (\deg_{C_2},\deg_{C_6},\deg_{C_8},\deg_l)
                                      =(1,2,1,6).             \tag{S-008.5}
\]
The portion of the quadratic coefficient independent of $C_8$ has
weight ten; on this slice its only possible multidegrees are
$C_2^5$ and $C_2^2C_6$.  Neither can contribute to (S-008.5),
whose $C_2$-degree is one.  The linear nonic coefficient has weight
19, which cannot be formed from the even weights $2,6,8$, so it
vanishes on this slice.  The constant term is independent of $l$.
Thus (S-008.5) isolates the residual symbol evaluated on two copies
of $l^3C_6$ and one copy of $C_8$, together with the free $C_2$.
The nonzero factor $\binom96^2=84^2$ from the padded formula does
not change injectivity.

The highest symbol of the exceptional summand is
\[
                         \Delta^7E(a_0G+b_0F).
\]
Put $l=e_0$ and $a_2=b_2=c_0=0$.  This becomes
\[
                         2E^8a_0b_0c_1c_2^7.
\]
Its coefficient at
$a_0^6a_1^3b_0^4b_1^5c_1c_2^7$ is
$-2\binom83=-112$.  Both nonic letters have zeroth exponent at
least three.  Symmetric polarization therefore proves a nonzero
equivariant restriction from the irreducible $S_{(10,9,7)}$ to
functions on $(l,C_6,C_8)$, hence an injective restriction.
Tensoring this injection with the identity on the independent
$C_2\in\operatorname{Sym}^2E$ kills every remaining copy of $L$.
The residual quadratic coefficient is independent of $C_8$, so
(S-008.3) eliminates it and the affine-linear proof eliminates
the remainder.  This proves the refinement, and together with
QC-010 gives $q=a$ throughout $\beta_1\leq11$, $|\beta|\leq28$.
The result for $(12,8,8)$ is the separate OBS-006 statement.

\subsection{S-009. A version-scoped correction to a sufficient padding bound}
\paragraph{Inspected source and precise correction.}
Kadish--Landsberg, \emph{Padded polynomials, their cousins, and
geometric complexity theory}, arXiv:1204.4693v1, Theorem~1.3,
printed p.~2, states a sufficient no-ideal-copy condition
\[
                         p_1\geq\min\{d(n-1),dn-m\}.          \tag{S-009.1}
\]
The proof's use of stability for the first branch is on printed
p.~5, in the paragraph beginning ``The multiplicity of $S_\pi W$''.
The branch $p_1\geq d(n-1)$ is false as stated.  This correction
is expressly about the inspected arXiv version.  The publisher
full text returned HTTP~403, so no statement about the final
published wording is made.

\paragraph{Complete two-copy counterexample.}
Take $n=3$, $m=1$, $d=6$, and the binary cubic with ordinary
coefficients
\[
                       f=ax^3+bx^2y+cxy^2+dy^3.
\]
To avoid confusion between the outer degree and the final coefficient,
the latter $d$ is treated as a formal variable in the following
identities.  Set
\[
 H=b^2-3ac,\qquad
 \Delta=b^2c^2-4ac^3-4b^3d-27a^2d^2+18abcd.
\]
The raising derivation in these ordinary coefficient coordinates is
\[
                         D=3a\partial_b+2b\partial_c+c\partial_d.
\]
Direct differentiation gives $DH=D\Delta=0$.  Their torus weights
are $(4,2)$ and $(6,6)$, respectively.  Hence
\[
                              P=a^2\Delta,\qquad Q=H^3        \tag{S-009.2}
\]
are highest-weight polynomials of coefficient degree six and weight
$(12,6)$.  They are independent: $P(x^3+y^3)=-27$, while
$P(x^2y)=0$ and $Q(x^2y)=1$.
The ambient multiplicity is exactly
\[
 a_{(12,6)}(6[3])=p_{\leq6,\leq3}(6)-p_{\leq6,\leq3}(5)
                         =7-5=2.                             \tag{S-009.3}
\]
Here the two partition counts are the binary weight multiplicities,
so subtracting consecutive weights is the exact $\mathrm{GL}_2$
highest-weight calculation.

Every form $l^2q_1$ has a repeated linear factor, and substituting
$l=ux+vy$, $q_1=sx+ty$ into $\Delta$ gives zero identically.
Thus $P$ is a nonzero ideal vector of the padded variety.  The
form $x^2y$ lies in that variety and $Q$ is nonzero there.  The
entire two-dimensional space is therefore controlled:
\[
                a_{(12,6)}=2,\qquad q_{(12,6)}=1,\qquad
                \operatorname{mult}_{(12,6)}I(V_{1,3})_6=1.   \tag{S-009.4}
\]
This argument also works in any larger ambient vector space by
restriction to the first two coordinates.  But the sufficient
condition in (S-009.1) holds, since
$12=\min\{6(3-1),18-1\}$.  Consequently (S-009.4) is a complete
counterexample, not merely a dimension-based objection to its proof.
The standalone \texttt{check\_kl\_counterexample.py} checks all
raising, substitution, weight and partition-count identities over
the integers.  The root independently performed the same verification
in exact BigInt polynomial arithmetic at 2026-11-22 01:08 UTC.

\paragraph{What is retained.}
This counterexample does not affect the necessary condition
$p_1\geq d(n-m)$ for a copy to occur in the padded coordinate ring,
or the associated essential-variable bound (with S-004's external
padding convention).  It does not affect the padding multiplicity
map separately proved in GCT\_KRON O-012.b: that proof exhibits an
injective highest-weight map and checks its invertible factorial
diagonal transpose.  O-012.c has its own flag-minor proof of its
stabilization range.  Thus O-012 and its F-014/F-015 consequences
remain valid; their citation of Theorem~1.3 must refer only to the
separately valid padding inequality, not to (S-009.1).

The other sufficient branch, $p_1\geq dn-m$, says that the total
tail size is at most $m$ and follows directly from coefficient
support and free lower jets.  It is retained.  The supplied BIP
Theorem \texttt{pro:I-KaLa} and IP Theorem~1.8 use only the
necessary shape condition and remain intact.  No equality between
the actual coordinate ring and a normalization or all-global-section
target is inferred.  The detailed search and scope classification are
recorded in \texttt{kl\_dependency\_audit.tsv}.

The later 2016 author book prepublication version
\emph{Geometry and Complexity Theory}, Section~8.4.2,
Propositions~8.4.2.1 and 8.4.2.4, printed pp.~235--236, states the
necessary conditions and the external-variable count $m^2+1$;
the inspected passage does not include the false sufficient branch.
Exact primary URLs and version locators appear in the source audit.

% Complete certificate; all graph, entry, packed-rank and scalar-rank checks passed.
\subsection*{S-010. A complete equal-tail-nine certificate}
\paragraph{Statement.}
For
every admissible $d\geq4$ and $\lambda_d=(9d-27,9,9,9)$,
\[
                         K_9(\lambda_d)=q_{\lambda_d}=a_{\lambda_d}.
\]
The common multiplicity is 2386 for $d\geq16$.  The defining equation
of the first-eight-jet determinant image from S-005 then has determinant
character $c_G\geq10$.  All exact graph, entry and rank checks described
below have passed.

\paragraph{The complete centered space.}
Put $E=\mathbb C^3$ and
\[
 R=\operatorname{Sym}\left(\bigoplus_{j=2}^9\operatorname{Sym}^jE\right).
\]
The $\det^9$ relative-invariant space has weighted degree27, where
the weight of $C_j$ is $j$.  The exact multigraded six-term Weyl
calculation \texttt{invariant\_profiles.py --tail 9} gives dimension
2386 in 230 nonzero profiles $(k_2,\ldots,k_9)$; its largest block
has dimension93.  The independent direct $\mathrm{GL}_4$ character
calculation in \texttt{check\_equal\_tail\_stabilization.py} gives
the same ambient multiplicity at outer degree27, and proves that
degree16 is the first stable degree.

Each supplied polynomial is defined by a graph with nine alternating
triples.  A vertex of valency $j$ represents the symmetric tensor
$C_j$, and no triple contains the same vertex twice.  Sum the product
of the tensor entries over all colors assigned to the slots, with
the product of the nine alternating symbols as sign.  This defines
an integral polynomial $I_\Gamma$ in the full, unspecialized centered
jets.  The alternating-symbol identity shows directly that it has
$\mathrm{GL}_3$ character $\det^9$ and is killed by the simple raising
operators.  No assertion that an arbitrary restricted polynomial is
highest-weight is used.

\paragraph{An injective quadratic normalization.}
To store and construct these polynomials efficiently, specialize the
symmetric matrix of $C_2$ to $I_3$.  Restriction to this slice is
injective on the fixed $\det^9$ space.  Indeed every nonsingular
complex symmetric quadratic form is congruent to $I_3$; if a relative
invariant vanished on the slice, covariance would make it vanish for
every nonsingular quadratic form and hence identically by density.
The odd determinant character causes no problem: an orthogonal
reflection changes the other tensors as well and multiplies the
relative invariant by $-1$.

At $C_2=I_3$, each quadratic vertex forces its two slot colors to
coincide.  The actual graph, including all quadratic vertices, remains
part of the certificate.  Only the restricted coefficient table omits
those quadratic factors.  Thus the original multidegree and its
polynomial lift are not lost under normalization.

\paragraph{Signed color-orbit compression and its exact check.}
Fix one alternating triple to colors $(0,1,2)$ and expand the remaining
eight triples.  Because there are nine triples, the global color
permutation character is $\operatorname{sgn}\sigma$.  For a monomial
orbit $O$ with no odd permutation in its stabilizer, let $S_O$ be the
sum of the fixed-triple coefficients, each transported to a canonical
representative with the sign of that transport.  The coefficient of
the representative in the full restricted polynomial is
\[
                              \frac6{|O|}S_O.
\]
All other coefficients follow by the permutation sign.  If the
stabilizer contains an odd permutation the full coefficient is zero,
and the orbit is discarded.  This rule follows by restoring the six
possible orientations of the fixed triple; the number of maps sending
any orbit element to the representative is $6/|O|$, all with the same
sign when the stabilizer is even.

The standard-library checker \texttt{verify\_bracket\_basis\_c9.py}
independently reconstructs every graph with exact integer arithmetic,
enforcing the quadratic equal-color constraints and the signed orbit
rule.  It compares every coefficient with the saved tables and checks
the nonzero coefficient minor in every profile modulo101.  This is
an independent expansion from the actual nine-triple graphs, rather
than a raising check on specialized coefficient tables.  The latter
alone would not establish the original highest-weight property.

The nonzero restricted coefficient minors prove linear independence
of the original graph polynomials over $\mathbb Q$.  Their total
number equals the complete ambient bound2386, so they form a basis.
For a graph with $k=\sum_j k_j$ vertices, its explicit polynomial
lift is
\[
                   L(I_\Gamma)=F_0^{27-k}I_\Gamma(C_2,\ldots,C_9).
\]
Each $C_j$ has denominator dividing $F_0^{j-1}$, so the displayed
expression is polynomial, of coefficient degree27 and highest weight
$(216,9,9,9)$.  S-005's centering argument and the matching direct
character count prove that all2386 lifts form the entire multiplicity
space in that degree.

\paragraph{Determinant points with the required residue-field normalization.}
The deterministic dataset consists of2420 integer pencils
\[
                      \det(x_0I_9+x_1A_1+x_2A_2+x_3A_3).
\]
Every $A_i$ has exactly zero integer trace, hence $F_0=1$ and
$F_1=0$ over $\mathbb Q$.  The normalized second tensor is
$-\operatorname{tr}(A(v)^2)/72$.  Over $\mathbb F_{101}$ the sampler
chooses a nonsingular trace Gram matrix of square determinant,
performs exact congruence to $I_3$, and applies the resulting variable
change to the actual matrices.  Replacing the last diagonal entry
by the negative sum of the other eight preserves the residue class
and gives exact integer trace zero.

Thus the actual integer lifts need only satisfy $C_2\equiv I_3$
modulo101, which they do.  There is no assertion that the lifted
quadratic forms equal $I_3$ over $\mathbb Q$.  The original graph
polynomials evaluated at these lifts reduce to their stored restricted
tables.  All tensor normalizations have denominators dividing $9!$,
which is invertible modulo101.

The complete standard-library sampler was rerun independently by this
agent: all7260 direct determinant checks passed, and all bytes matched
the supplied2420-point dataset.  Its normalized tensor SHA-256 is
\begin{center}\small\ttfamily
f30ffaec897ba3ddca20d9a7944477027f01f49da97fe0d5df307d81f0c24d82.
\end{center}
The full serialized dataset SHA-256 is
\begin{center}\small\ttfamily
25781148d5977ab0d07d45b7163c74fc435e43a3530945d40a320444dd6179e8.
\end{center}
The exact congruence, trace formula, interpolation and integer-lift
conventions are in \texttt{generate\_orthonormal\_samples.py} and
QC-008; the rerun comparison is saved separately.

\paragraph{Final rank gate.}
The full matrix has2386 rows and2420 columns, rank2386 over
$\mathbb F_{101}$, and selected maximal minor14.  The first2386
columns already give that nonzero minor.  Its row-major byte SHA-256 is
\begin{center}\small\ttfamily
d781202425cd3e9bba02d996afa76de3b1bda036564bab797f3ac7135d439b50.
\end{center}
The signed standard-library entry checker independently restored all
six color images and recomputed every one of the5,774,120 matrix
entries; every entry matched.  A separate ordinary scalar Gaussian
calculation gave the same rank2386 and minor14, independently of the
packed calculation.  The exact standard-library graph reconstruction
also completed all2386 polynomials and all230 coefficient minors.
The packed rank took125.188seconds, and the independent scalar rank
383.717seconds.  Per-profile source hashes bind the completed graph
checks and any resumed prefixes to their exact input tables.

Every pencil is a linear endomorphism specialization of $\det_9$ and
belongs to its orbit closure by density of invertible endomorphisms.
The nonzero modular minor lifts to a nonzero rational
minor on genuine determinant pencils.  Hence no nonzero highest-weight
vector of the indicated type vanishes on the determinant orbit closure.
S-007 gives padded equality for the whole second-row-at-most-nine
region, so the degree27 comparison is $K=q=a=2386$.

Multiplication by $F_0$ is injective both in the polynomial ring and
the integral determinant coordinate ring.  The complete centered upper
bound is2386, and the actual multiplicity already reaches it in
degree16.  A nonzero ideal vector in an earlier degree would remain
nonzero after multiplication to degree27; in every later degree the
complete highest-weight space is the corresponding $F_0$ multiple.
Thus the zero kernel in degree27 implies the stated equality in every
admissible degree.  Finally, a first-eight-jet hypersurface equation
with character9 would give a nonzero element of this very centered
space, independent of $C_9$, vanishing on determinants.  Full rank
excludes it; S-006 already excluded smaller characters.  Therefore
the hypersurface character is at least10.

\paragraph{Scope.}
This is a calculation of the actual homogeneous ideal and restriction
map, not a comparison with a normalization or orbit-coordinate-ring
upper bound.  The graph covariance is a polynomial Schur-functor
construction, while the localization and saturation concern coherent
graded ideals.  The SGA~3 I4.7.2--4.7.3 and cote133, scan028
dictionaries are used in the same precise scope as S-005--S-007.
The modular certificate establishes the characteristic-zero assertion;
it does not assert all-characteristic equality or any motivic statement.


\subsection{S-011. The degree-six nonic multiplicity-four gate}
\paragraph{Proved statement.}
For $n=9$, $d=6$, and $\lambda=(15,15,12,12)$,
\[
                    K_9(\lambda)=q_\lambda=a_\lambda=4.
\tag{S-011.1}
\]
Thus this label is not an actual multiplicity obstruction at $(8,9)$.
The calculation below treats all four copies explicitly.

\paragraph{Four integral highest-weight polynomials.}
Use six vector letters $y_0,\ldots,y_5\in\mathbb C^4$.
A four-bracket $[ijkl]$ is the determinant of these four letters,
in the written order. A two-bracket $[ij]$ is the determinant of
their first two coordinates. Put
\[
\begin{split}
\Omega_1={}&[0123]^4[0145]^4[0234][2345]^3
                         [12][35][45],\\
\Omega_2={}&[0123]^4[0124][0145]^3[0234][2345]^3
                         [15][35][45],\\
\Omega_3={}&[0123]^3[0145]^3[0234]^3[1245][2345]^2
                         [15]^2[35],\\
\Omega_4={}&[0123]^3[0124]^2[0145]^3[0234][1235][2345]^2
                         [35]^2[45].
\end{split}                                                     \tag{S-011.2}
\]
Every letter has degree9, and every expression has twelve
four-brackets and three two-brackets. Define
\[
 P_i(F)=
 \left.
 \Omega_i(\partial_{x^{(0)}},\ldots,\partial_{x^{(5)}})
             \prod_{r=0}^5 F(x^{(r)})
 \right|_{x^{(0)}=\cdots=x^{(5)}=0}.
\tag{S-011.3}
\]
The differential pairing is equivariant. A top two-minor and a full
four-minor are annihilated by all simple raising operators, so each
$P_i$ is either zero or a highest-weight vector of weight
$(15,15,12,12)$ in coefficient degree6. Formula (S-011.3) fixes the
normalization: in ordinary coefficient coordinates a term
$\prod_r y_r^{\alpha_r}$ contributes
$\prod_r\alpha_r![x^{\alpha_r}]F$.
Equivalently, using normalized symmetric-tensor coordinates, it is
the common scalar $(9!)^6$ times the usual symmetrized symbol.
The nonzero restriction minor below proves that all four $P_i$
are nonzero and independent.

\paragraph{Three coefficient functionals on a determinant subfamily.}
The checked ternary determinant theorem (CUE\_GCT G-05,
GCT\_KRON O-007; Beauville, Proposition3.1(a)) puts every ternary
septic in the degree-seven determinant orbit closure.
Multiplication by $x_2^2$, using the block diagonal extension with
two $1$ by $1$ blocks $x_2$, therefore puts
\[
                      x_2^2G(x_0,x_1,x_3)
\tag{S-011.4}
\]
in $\mathcal D_{9,4}$ for every ternary septic $G$.
This uses closedness and continuity, so no choice of a determinantal
representation for an individual singular septic is required.

For $j=1,2,3$, let the six integer vectors
$v_{j0},\ldots,v_{j5}\in\mathbb Z^3$ be the rows in the following arrays:
\[
\begin{split}
V_1={}&((93,13,44),(97,40,38),(63,76,74),
         (85,14,99),(52,42,35),(25,12,10)),\\
V_2={}&((89,42,96),(77,40,46),(78,54,1),
         (29,20,86),(86,98,95),(90,42,20)),\\
V_3={}&((89,65,6),(69,75,100),(96,67,68),
         (87,74,78),(33,69,60),(100,32,27)).
\end{split}
\]
Define the polynomial families and linear functionals
\[
\begin{split}
 F_j(t;x)&=x_2^2\sum_{r=0}^5
       t_r\bigl(v_{jr}\cdot(x_0,x_1,x_3)\bigr)^7,\\
 L_j(P)&=[t_0t_1t_2t_3t_4t_5]\,P(F_j(t)).
\end{split}                                                     \tag{S-011.5}
\]
Every member of every family lies in $\mathcal D_{9,4}$ by
(S-011.4). Consequently an equation of that variety is annihilated
by each $L_j$.

Here is a short exact contraction for (S-011.5).
For an ordered assignment of the six vectors to the six letters,
substitute
\[
          y_r=(v_{r0},v_{r1},z_r,v_{r2})
\]
in $\Omega_i$, and take the coefficient of
$z_0^2z_1^2\cdots z_5^2$. Each four-bracket is linear in the $z_r$:
the coefficient of its letter in position $a$ is
$(-1)^{a+2}$ times the corresponding ternary determinant, with
positions numbered from zero. The two-brackets do not involve $z$.
A dynamic program multiplying the twelve linear factors and
discarding exponents greater than two has only $3^6=729$ states.
Sum this coefficient over all $6!=720$ ordered assignments.
The differential normalization then multiplies the sum by
\[
                         (2!\,7!)^6=37\pmod {101}.             \tag{S-011.6}
\]
The resulting unscaled sums, in rows $i$ and columns $j$, are
\[
 C=
 \begin{pmatrix}
 20&15&81\\
 48&59&27\\
 37&100&91\\
 100&12&17
 \end{pmatrix}\pmod {101}.
\tag{S-011.7}
\]
Thus the actual coefficient-functional block is $37C$.
The division implicit in reporting $C$ is only a presentation
normalization; the final matrix below uses the literal functionals
$L_j$ of (S-011.5).

\paragraph{One explicit determinant evaluation.}
Put $X=x_0x_1x_2$ and
\[
 F=X^2(x_3^3+X),\qquad
 M=\begin{pmatrix}x_3&x_0&0\\0&x_3&x_1\\x_2&0&x_3\end{pmatrix}.
\]
Since $\det M=x_3^3+X$, the form $F$ is the determinant of the
block matrix with six scalar diagonal blocks
$x_0,x_0,x_1,x_1,x_2,x_2$ and the final block $M$.
Let
\[
 g=\begin{pmatrix}
 1&1&1&0\\
 0&0&0&1\\
 0&1&0&0\\
 0&0&1&0
 \end{pmatrix},\qquad \det g=1,
\]
and use the actual integer determinant
\[
 \widehat F(x)=F(g^{\mathsf T}x)
  =\bigl(x_0(x_0+x_2)(x_0+x_3)\bigr)^2
    \bigl(x_1^3+x_0(x_0+x_2)(x_0+x_3)\bigr).                 \tag{S-011.8}
\]
By the chain rule, $P_i(\widehat F)$ is obtained from the original
$F$ by replacing each two-bracket with the bracket of the two
linear coordinates $y_0+y_1+y_2$ and $y_3$; the full four-brackets
are unchanged because $\det g=1$.

Write the two ordinary monomials of $F$ as
$A=x_0^2x_1^2x_2^2x_3^3$ and
$E=x_0^3x_1^3x_2^3$. Every changed two-bracket contributes exactly
one $y_3$ coordinate. Together with the twelve full brackets this
requires total fourth-coordinate weight15, so among six selected
monomials only $A^5E$ can contribute. The required extra color
counts from the three two-brackets are $(1,1,1,3)$.
Exact coefficient joins give sums
\[
                           (2,81,30,96)\pmod {101}
\]
over the six placements of $E$. The common ordinary differential
factor is
\[
 (2!2!2!3!)^5(3!3!3!0!)=91\pmod {101},
\]
and therefore
\[
                  (P_1(\widehat F),\ldots,P_4(\widehat F))
                       =(81,99,3,50)\pmod {101}.             \tag{S-011.9}
\]
The implementation uses only positive letter automorphisms when
grouping placements. Coordinate permutations, whose signs can be
negative for three two-brackets, are not silently identified.

\paragraph{Complete rank and ambient count.}
Combining the three literal coefficient functionals with the
one literal evaluation gives
\[
 \bigl(L_1(P_i),L_2(P_i),L_3(P_i),P_i(\widehat F)\bigr)_{i=1}^4
 =
 \begin{pmatrix}
 33&50&68&81\\
 59&62&90&99\\
 56&64&34&3\\
 64&40&23&50
 \end{pmatrix}
 \pmod {101},
 \qquad \det=21\ne0.                                        \tag{S-011.10}
\]
All entries before reduction are integers. Thus the same minor is
nonzero in characteristic zero, and all four highest-weight
polynomials remain independent on the actual determinant closure.

Independently, exact monomial counting in
$\operatorname{Sym}^6(\operatorname{Sym}^9\mathbb C^4)$ followed by
the $24$-term Weyl alternation gives
\[
 \sum_{\sigma\in S_4}\operatorname{sgn}(\sigma)
   \dim W_{\lambda+\rho-\sigma\rho}=4,\qquad \rho=(3,2,1,0).
\]
The complete24 integer weight counts are saved in
\texttt{four\_nonic\_ambient\_character.json}; no floating arithmetic
is used. The restriction rank is therefore exactly the ambient
multiplicity4. QC-011 identifies $q=a$ in coefficient degrees at
most9 for the actual $(8,9)$ padded closure. In particular $q=4$
in degree6, proving (S-011.1).

\paragraph{Reproducibility and scope.}
The selected symbols and exact ambient-selection data are in
\texttt{four\_nonic\_ambient\_selection.json}.
The ternary coefficient checker is
\texttt{check\_nonic\_a4\_ternary\_slice.py};
the sparse determinant contraction is
\texttt{evaluate\_four\_nonic\_sparse.py}, using the readable
standard-library determinant-power joins in the quaternary branch.
The final complete checker is
\texttt{verify\_four\_nonic\_certificate.py}.
It reconstructs the integer symbols, all three coefficient columns,
the sparse determinant column, the exact character upper bound,
and the final minor, with no NumPy or black-box representation package.
The intermediate rank3 on twelve tested square-times-ternary
polarizations was only a diagnostic; the proof uses the nonzero
four-column minor (S-011.10), not an unproved upper bound on that
smaller family's image.
